\clearpage
\SourceLeaf{264}
\thispagestyle{empty}
\pdfbookmark[0]{Untersuchungen über Themen der höheren Geodäsie. Erste Abhandlung}{geodaesie1}
\vspace*{.3in}
\begin{center}
{\large UNTERSUCHUNGEN}\par\bigskip
ÜBER\par\bigskip
{\Large THEMEN DER HÖHEREN GEODÄSIE}\par\vspace{25pt}
ERSTE ABHANDLUNG\par\vspace{35pt}
VON\par\vspace{35pt}
{\large CARL FRIEDRICH GAUSS}\par\vspace{35pt}
{\small DER KÖNIGLICHEN GESELLSCHAFT AM 23. OKTOBER 1843 VORGELEGT.}
\end{center}
\vfill
\begin{center}
\rule{.8\linewidth}{.3pt}\par\smallskip
Abhandlungen der Königl. Gesellschaft der Wissenschaften zu Göttingen. Band II.\par
Göttingen 1844.\par\smallskip
\rule{.8\linewidth}{.3pt}
\end{center}
\par\bigskip\begingroup\fontsize{9}{11}\selectfont\raggedright\noindent Die erste Abhandlung zur höheren Geodäsie ist aus \textit{Werke}, Band IV, in modernes Deutsch übertragen. Die beigefügten Tafeln sind derselben Ausgabe entnommen. Das historische Original bleibt gesondert erhalten. Der Anfang des \href{https://doi.org/10.3931/e-rara-2860}{Separatdrucks von 1844 der ETH-Bibliothek} wurde im OCR-Text verglichen. Terminologisch wurden \href{https://arxiv.org/abs/math/0409578}{Erhard Scholz} (2004), \href{https://doi.org/10.12902/zfv-0387-2022}{Miljenko Lapaine} (2022) und \href{https://en.wikisource.org/w/index.php?oldid=12858814}{J. W. L. Glaishers Beitrag zu mathematischen Tafeln} (1911) herangezogen.\par\endgroup
\clearpage
